\section{CSLib and its named representation of the lambda calculus}
\label{sec:cslib}

\subsection{CSLib}
\label{subsec:cslib-overview}

CSLib is a community-maintained library of foundational computer
science formalizations for the Lean~4 theorem prover, built on top
of Mathlib. \cite{demoura2026a} 
Its goal is to provide idiomatic, well-reviewed, and actively
maintained formalizations of standard results, so that later
research projects can build directly on a shared, trusted base
rather than re-deriving basic definitions and lemmas from scratch.
The work reported in this paper concerns CSLib's treatment of the
untyped lambda calculus, and in particular its representation(s) of
variable binding.

\subsection{Named terms}
\label{subsec:named-terms}

CSLib currently represents untyped lambda terms in a standard named
style: variables are drawn from a countably infinite set of atoms,
and terms are built from variables, application, and abstraction. TODO dummy math

\begin{definition}[Named lambda terms]
\label{def:named-terms}
Let $\mathrm{Var}$ be a countably infinite set of variables. The set
$\Lam$ of named lambda terms is generated by
\[
  t \;::=\; x \mid t\, t \mid \lambda x.\, t
  \qquad (x \in \mathrm{Var}).
\]
\end{definition}

On top of this definition, CSLib provides the usual auxiliary
notions needed to state and prove properties of the calculus: free
variables, capture-avoiding substitution, and a baseline notion of
alpha equivalence used as the reference point against which
Crole's alternative definitions are shown equivalent in
Section~\ref{sec:alpha-equivalences}.

